% ReCLamO-Harness paper -- DRAFT. Not for submission.
% Every result number lives in tables/*.tex, generated by analysis/make_tables.py.
% Main-study cells are \PENDING until `make_tables.py --stage main` has been run on
% the completed grid; pilot numbers (seeds 2-3, n=16 per cell) are labelled preliminary.
\documentclass[11pt]{article}

\usepackage[margin=1in]{geometry}
\usepackage[utf8]{inputenc}
\usepackage[T1]{fontenc}
\usepackage{lmodern}
\usepackage{amsmath, amssymb}
\usepackage{setspace}
\usepackage{microtype}
\usepackage{enumitem}
\usepackage{authblk}
\usepackage{booktabs}
\usepackage{tabularx}
\usepackage{xcolor}
\usepackage{graphicx}
\usepackage[numbers,sort&compress]{natbib}
\usepackage{hyperref}

\setstretch{1.15}

\hypersetup{
  colorlinks=true,
  linkcolor=black,
  citecolor=black,
  urlcolor=blue
}

% ---------------------------------------------------------------- draft markers
\newcommand{\PENDING}[1]{\textcolor{orange!80!black}{\textbf{[PENDING:} #1\textbf{]}}}
\newcommand{\TODO}[1]{\textcolor{red!70!black}{\textbf{[TODO:} #1\textbf{]}}}
\newcommand{\pendingtable}[2]{%
  \begin{table}[t]\centering
  \fbox{\parbox{0.9\linewidth}{\centering\vspace{1ex}\PENDING{#2}\vspace{1ex}}}
  \caption{#1}
  \end{table}}

% Generated numbers. Pilot macros always exist; main macros appear once the grid is done.
% Generated by paper/analysis/make_tables.py --stage pilot. Do not edit.
\newcommand{\pilotExactplainThreeTwoksmall}{9/16}
\newcommand{\pilotWilsonplainThreeTwoksmall}{[0.33, 0.77]}
\newcommand{\pilotMeanplainThreeTwoksmall}{0.60}
\newcommand{\pilotSubplainThreeTwoksmall}{0.00}
\newcommand{\pilotDnfplainThreeTwoksmall}{2}
\newcommand{\pilotExactplainThreeTwokmedium}{0/16}
\newcommand{\pilotWilsonplainThreeTwokmedium}{[0.00, 0.19]}
\newcommand{\pilotMeanplainThreeTwokmedium}{0.00}
\newcommand{\pilotSubplainThreeTwokmedium}{--}
\newcommand{\pilotDnfplainThreeTwokmedium}{16}
\newcommand{\pilotExactplainOneThreeOneksmall}{8/16}
\newcommand{\pilotWilsonplainOneThreeOneksmall}{[0.28, 0.72]}
\newcommand{\pilotMeanplainOneThreeOneksmall}{0.56}
\newcommand{\pilotSubplainOneThreeOneksmall}{0.00}
\newcommand{\pilotDnfplainOneThreeOneksmall}{0}
\newcommand{\pilotExactplainOneThreeOnekmedium}{7/16}
\newcommand{\pilotWilsonplainOneThreeOnekmedium}{[0.23, 0.67]}
\newcommand{\pilotMeanplainOneThreeOnekmedium}{0.47}
\newcommand{\pilotSubplainOneThreeOnekmedium}{0.00}
\newcommand{\pilotDnfplainOneThreeOnekmedium}{2}
\newcommand{\pilotExacthvZeroOneThreeTwoksmall}{11/16}
\newcommand{\pilotWilsonhvZeroOneThreeTwoksmall}{[0.44, 0.86]}
\newcommand{\pilotMeanhvZeroOneThreeTwoksmall}{0.71}
\newcommand{\pilotSubhvZeroOneThreeTwoksmall}{0.75}
\newcommand{\pilotDnfhvZeroOneThreeTwoksmall}{0}
\newcommand{\pilotExacthvZeroOneThreeTwokmedium}{9/16}
\newcommand{\pilotWilsonhvZeroOneThreeTwokmedium}{[0.33, 0.77]}
\newcommand{\pilotMeanhvZeroOneThreeTwokmedium}{0.62}
\newcommand{\pilotSubhvZeroOneThreeTwokmedium}{0.38}
\newcommand{\pilotDnfhvZeroOneThreeTwokmedium}{0}
\newcommand{\pilotExacthvZeroTwoThreeTwoksmall}{10/16}
\newcommand{\pilotWilsonhvZeroTwoThreeTwoksmall}{[0.39, 0.82]}
\newcommand{\pilotMeanhvZeroTwoThreeTwoksmall}{0.68}
\newcommand{\pilotSubhvZeroTwoThreeTwoksmall}{1.62}
\newcommand{\pilotDnfhvZeroTwoThreeTwoksmall}{0}
\newcommand{\pilotExacthvZeroTwoThreeTwokmedium}{7/16}
\newcommand{\pilotWilsonhvZeroTwoThreeTwokmedium}{[0.23, 0.67]}
\newcommand{\pilotMeanhvZeroTwoThreeTwokmedium}{0.49}
\newcommand{\pilotSubhvZeroTwoThreeTwokmedium}{1.44}
\newcommand{\pilotDnfhvZeroTwoThreeTwokmedium}{0}
\newcommand{\pilotExacthvZeroTwoOneThreeOneksmall}{10/16}
\newcommand{\pilotWilsonhvZeroTwoOneThreeOneksmall}{[0.39, 0.82]}
\newcommand{\pilotMeanhvZeroTwoOneThreeOneksmall}{0.69}
\newcommand{\pilotSubhvZeroTwoOneThreeOneksmall}{2.44}
\newcommand{\pilotDnfhvZeroTwoOneThreeOneksmall}{0}
\newcommand{\pilotExacthvZeroTwoOneThreeOnekmedium}{8/16}
\newcommand{\pilotWilsonhvZeroTwoOneThreeOnekmedium}{[0.28, 0.72]}
\newcommand{\pilotMeanhvZeroTwoOneThreeOnekmedium}{0.57}
\newcommand{\pilotSubhvZeroTwoOneThreeOnekmedium}{2.06}
\newcommand{\pilotDnfhvZeroTwoOneThreeOnekmedium}{0}
\newcommand{\pilotPhvZeroOneThreeTwokvsplainOneThreeOneksmall}{0.45}
\newcommand{\pilotBChvZeroOneThreeTwokvsplainOneThreeOneksmall}{5/2}
\newcommand{\pilotPhvZeroOneThreeTwokvsplainOneThreeOnekmedium}{0.50}
\newcommand{\pilotBChvZeroOneThreeTwokvsplainOneThreeOnekmedium}{2/0}
\newcommand{\pilotPhvZeroOneThreeTwokvsplainOneThreeOnekpooled}{0.18}
\newcommand{\pilotBChvZeroOneThreeTwokvsplainOneThreeOnekpooled}{7/2}
\newcommand{\pilotPhvZeroOneThreeTwokvsplainThreeTwoksmall}{0.73}
\newcommand{\pilotBChvZeroOneThreeTwokvsplainThreeTwoksmall}{5/3}
\newcommand{\pilotPhvZeroOneThreeTwokvsplainThreeTwokmedium}{0.004}
\newcommand{\pilotBChvZeroOneThreeTwokvsplainThreeTwokmedium}{9/0}
\newcommand{\pilotPhvZeroOneThreeTwokvsplainThreeTwokpooled}{0.01}
\newcommand{\pilotBChvZeroOneThreeTwokvsplainThreeTwokpooled}{14/3}
\newcommand{\pilotPhvZeroTwoOneThreeOnekvsplainOneThreeOneksmall}{0.50}
\newcommand{\pilotBChvZeroTwoOneThreeOnekvsplainOneThreeOneksmall}{2/0}
\newcommand{\pilotPhvZeroTwoOneThreeOnekvsplainOneThreeOnekmedium}{1.00}
\newcommand{\pilotBChvZeroTwoOneThreeOnekvsplainOneThreeOnekmedium}{3/2}
\newcommand{\pilotPhvZeroTwoOneThreeOnekvsplainOneThreeOnekpooled}{0.45}
\newcommand{\pilotBChvZeroTwoOneThreeOnekvsplainOneThreeOnekpooled}{5/2}
\newcommand{\pilotPhvZeroOneThreeTwokvshvZeroTwoOneThreeOneksmall}{1.00}
\newcommand{\pilotBChvZeroOneThreeTwokvshvZeroTwoOneThreeOneksmall}{3/2}
\newcommand{\pilotPhvZeroOneThreeTwokvshvZeroTwoOneThreeOnekmedium}{1.00}
\newcommand{\pilotBChvZeroOneThreeTwokvshvZeroTwoOneThreeOnekmedium}{2/1}
\newcommand{\pilotPhvZeroOneThreeTwokvshvZeroTwoOneThreeOnekpooled}{0.73}
\newcommand{\pilotBChvZeroOneThreeTwokvshvZeroTwoOneThreeOnekpooled}{5/3}
\newcommand{\pilotPhvZeroOneThreeTwokvshvZeroTwoThreeTwoksmall}{1.00}
\newcommand{\pilotBChvZeroOneThreeTwokvshvZeroTwoThreeTwoksmall}{1/0}
\newcommand{\pilotPhvZeroOneThreeTwokvshvZeroTwoThreeTwokmedium}{0.62}
\newcommand{\pilotBChvZeroOneThreeTwokvshvZeroTwoThreeTwokmedium}{3/1}
\newcommand{\pilotPhvZeroOneThreeTwokvshvZeroTwoThreeTwokpooled}{0.38}
\newcommand{\pilotBChvZeroOneThreeTwokvshvZeroTwoThreeTwokpooled}{4/1}

\IfFileExists{tables/main_macros.tex}{\input{tables/main_macros.tex}}{}

% Readable aliases for the pilot numbers used in the text.
\newcommand{\pVoneS}{\pilotExacthvZeroOneThreeTwoksmall}
\newcommand{\pVoneM}{\pilotExacthvZeroOneThreeTwokmedium}
\newcommand{\pPlainLongS}{\pilotExactplainOneThreeOneksmall}
\newcommand{\pPlainLongM}{\pilotExactplainOneThreeOnekmedium}
\newcommand{\pPlainShortS}{\pilotExactplainThreeTwoksmall}
\newcommand{\pVtwoLongS}{\pilotExacthvZeroTwoOneThreeOneksmall}
\newcommand{\pVtwoLongM}{\pilotExacthvZeroTwoOneThreeOnekmedium}
\newcommand{\pVtwoShortS}{\pilotExacthvZeroTwoThreeTwoksmall}
\newcommand{\pVtwoShortM}{\pilotExacthvZeroTwoThreeTwokmedium}
\newcommand{\pPVoneVsPlainLong}{\pilotPhvZeroOneThreeTwokvsplainOneThreeOnekpooled}
\newcommand{\pBCVoneVsPlainLong}{\pilotBChvZeroOneThreeTwokvsplainOneThreeOnekpooled}
\newcommand{\pPVoneVsVtwoLong}{\pilotPhvZeroOneThreeTwokvshvZeroTwoOneThreeOnekpooled}
\newcommand{\pBCVoneVsVtwoLong}{\pilotBChvZeroOneThreeTwokvshvZeroTwoOneThreeOnekpooled}
\newcommand{\pPVtwoVsPlainLong}{\pilotPhvZeroTwoOneThreeOnekvsplainOneThreeOnekpooled}
\newcommand{\pBCVtwoVsPlainLong}{\pilotBChvZeroTwoOneThreeOnekvsplainOneThreeOnekpooled}
\newcommand{\pSubVoneS}{\pilotSubhvZeroOneThreeTwoksmall}
\newcommand{\pSubVtwoLongS}{\pilotSubhvZeroTwoOneThreeOneksmall}

\newcommand{\code}[1]{\texttt{#1}}

\title{\textbf{Recursive Language Models on Local Hardware:\\A Replication-and-Measurement Study}\\[0.5ex]
\large\textcolor{red!70!black}{DRAFT --- results pending; not for citation}}
\author{Aaron Kingsley Clark\thanks{ORCID \href{https://orcid.org/0000-0002-6831-2663}{0000-0002-6831-2663}. Correspondence: \href{mailto:aaron.clark3@eastern.edu}{aaron.clark3@eastern.edu}.}}
\affil{Applied Artificial Intelligence\\Eastern University, St.~Davids, PA, USA}
\date{Draft of October 2026}

\begin{document}

\maketitle

\begin{abstract}
Recursive Language Models (RLMs) place a long input in a programming environment
as a variable and let a language model inspect it with code and call itself on
pieces of it, instead of reading the whole input in one prompt. The method was
introduced and evaluated with frontier models. We ask whether it carries over to a
locally served, quantized open-weight model on shared, serial hardware. We
contribute (i) ReCLamO-Harness, an open-source, clean-room RLM harness built for
local serving, with per-role endpoints, hard sub-call caps, context compaction that
guarantees every request fits the window, a forced finish on every limit, and a
compatibility mode that runs the original authors' fine-tuned RLM-Qwen3-8B
checkpoint; (ii) an evaluation set of eight long-context task generators written,
blind to the harness, by eight non-Anthropic models, together with three rounds of
audits that found and repaired answer-key and grader defects, and a scorer battery
that runs in continuous integration; and (iii) a pre-registered evaluation plan
comparing the harness, under two prompt and limit configurations, with a single-call baseline at 32K-token and 128K-token windows, with one primary
paired comparison per input size, a power-based target of 320 paired questions, and
pre-registered interim looks. A pilot on 16 questions per cell taught us that two
measurement errors, selecting the baseline's better attempt with the answer key and
defective graders, moved results by as much as the effects under study; we report
those lessons separately and draw no conclusions from the pilot.
\TODO{Findings sentence, to be written after the main grid (Section~\ref{sec:results}) completes.}
\end{abstract}

\noindent\textbf{Keywords:} recursive language models; long-context reasoning;
evaluation methodology; local inference; reproducibility

% ===========================================================================
\section{Introduction}
\label{sec:intro}

Language models degrade as their inputs grow, and they cannot read inputs longer
than their context window at all. Recursive Language Models
\citep{zhang2025rlm} address both problems with one move: the input is stored as a
variable in a persistent Python REPL, and the model writes code to peek at it,
split it, search it and pass pieces of it to fresh language-model calls
(\code{llm\_query}). The model sees only the short outputs of its own code. The
original paper reports that an RLM built on GPT-5 handles inputs far beyond that
model's window and outperforms the base model and common scaffolds on several
long-context benchmarks at comparable cost, and that a small fine-tuned model,
RLM-Qwen3-8B, improves on its base model.

Those results use frontier models served by commercial APIs, with sub-calls
answered by capable models. Many practitioners instead run a quantized open-weight
model on hardware they own, one or two requests at a time. This paper asks a
narrow, practical question: \emph{on such a setup, does an RLM harness answer
long-context questions more accurately than simply giving the model the whole
document, and at what cost?} It is a replication-and-measurement study, not a new
method. The study is also a record of how easily such a comparison is mismeasured:
in a pilot, our own measurement errors moved the result as much as the method did.

We make four contributions, and defer a fifth until the data are in.

\begin{enumerate}[leftmargin=*]
\item \textbf{An open-source harness for local serving}
  (Section~\ref{sec:harness}). ReCLamO-Harness \citep{reclamo} is a clean-room
  implementation of the RLM loop after \code{rlm-minimal} and \code{rlm}
  \citep{rlmminimal,rlmrepo}, tuned for Qwen-family models on a serial server:
  per-role endpoints (a planner and a reader on different servers), hard sub-call
  caps enforced in code, compaction that guarantees each request fits the window,
  a forced finish on every limit, and an \code{upstream-rlm-v0} mode that speaks
  the original scaffold's format so the released RLM-Qwen3-8B checkpoint
  \citep{rlmqwen8b} can plan without retraining.
\item \textbf{An independently authored, audited evaluation set}
  (Section~\ref{sec:tasks}). Eight task generators were written by eight
  non-Anthropic models, each given only a short brief and no access to the
  harness. Running them exposed task and grader defects that the generators' own
  self-tests could not catch; three audit rounds repaired them, and a scorer
  battery of 2,111 probes now runs in CI.
\item \textbf{A pre-registered evaluation design} (Section~\ref{sec:design}): paired
  arms on identical questions, one primary comparison per size sized by a power
  calculation from the pilot's discordance rate, interim looks with a corrected
  threshold, completion reported separately from accuracy, and a phased plan of
  validity checks, ablations, oracle interventions and coverage tests.
\item \textbf{Measurement lessons from a pilot} (Section~\ref{sec:lessons}),
  reported as such: a baseline that was best-of-2 selected on the answer key, grader
  bugs that cost one arm more than another, a prompt patch copied from the paper
  that suppressed the very behaviour it was meant to temper, and conditions under
  which the paper's 8B planner could not be fairly judged.
\item \textbf{Main results} (Section~\ref{sec:results}): \PENDING{accuracy, completion,
  cost and latency on up to 320 paired questions per primary comparison and size, and
  the outcome of each pre-registered hypothesis.}
\end{enumerate}

% ===========================================================================
\section{Background and related work}
\label{sec:related}

\paragraph{Recursive Language Models.}
\citet{zhang2025rlm} define an RLM as an inference-time scaffold in which the
prompt is part of the environment rather than the network's input. In their
instantiation the root model writes REPL code over a \code{context} variable,
calls sub-models with \code{llm\_query}, and ends by marking a variable or string
as final (\code{FINAL}/\code{FINAL\_VAR}). They evaluate on S-NIAH, built after
the single-needle task of RULER \citep{hsieh2024ruler}; OOLONG and the derived
OOLONG-Pairs \citep{bertsch2025oolong}; BrowseComp-Plus
\citep{chen2025browsecompplus}; and code question answering from LongBench v2
\citep{bai2025longbenchv2}. They fine-tune Qwen3-8B \citep{yang2025qwen3} on
trajectories from RLM(Qwen3-Coder-480B) on LongBench Pro
\citep{chen2026longbenchpro} to obtain RLM-Qwen3-8B. Their appendix of negative
results is directly relevant here: small models without strong coding ability
struggled as RLMs, thinking models ran out of output tokens, sequential sub-calls
were slow, and Qwen3-Coder needed an added warning not to make ``thousands of LM
subcalls for basic tasks''. A later post by the same group argues that harnesses
can make the individual model calls inside them in-distribution, a form of
compositional generalization \citep{zhang2026harness}. The authors also maintain the
reference implementations \citep{rlmrepo,rlmminimal} that ours follows.

\paragraph{Long-context evaluation.}
\citet{goldman2025retrieval} argue that many long-context benchmarks reduce to
retrieval, and that difficulty should be characterised by how dispersed and how
voluminous the needed information is; the RLM paper uses the same lens to explain
why GPT-5 degrades faster on OOLONG and OOLONG-Pairs than on S-NIAH. Context
degradation within the window (``context rot'') is documented by
\citet{hong2025contextrot}. Latent-structure tasks beyond haystacks are studied by
\citet{vodrahalli2024michelangelo}. \TODO{Cite ``lost in the middle'' (Liu et al.)
and LongBench v1 only after verifying the references; neither is in the source
reference lists used for this draft.}

\paragraph{Baselines the paper compares with.}
The RLM paper compares against a retrieval agent using BM25
\citep{robertson2009bm25}, a CodeAct agent \citep{wang2024codeact}, and summary
or compaction agents in the line of ReSum \citep{wu2025resum}, context folding
\citep{sun2025contextfolding} and MemAgent \citep{yu2025memagent}. Explicit memory
hierarchies \citep{packer2024memgpt} and interactive reading
\citep{chen2023memorymaze} are related ways to manage context. Recursive
delegation without a symbolic handle to the input appears in THREAD
\citep{schroeder2025thread} and ReDel \citep{zhu2024redel}. This draft does not yet
include a retrieval or compaction baseline; both are planned
(Section~\ref{sec:future}).

% ===========================================================================
\section{The ReCLamO harness}
\label{sec:harness}

ReCLamO-Harness is about 5,300 lines of Python under Apache-2.0. It was written
from the paper's description and the MIT-licensed reference code; where text is
copied verbatim (the upstream prompt for the compatibility mode, and the upstream
orchestrator addendum), the MIT notice sits beside it.

\subsection{The loop}
The context (a string, a list, or a dictionary of files) is loaded into a
persistent REPL as \code{context}. The system prompt describes the context's
type, size and chunk lengths, the REPL tools and the ways to finish. Each root turn
is one model call. Fenced \code{repl} blocks in the reply are executed in order;
their outputs, truncated to a configured length, are appended to the history. Inside
the REPL the model can call \code{llm\_query(prompt)},
\code{llm\_query\_batched(prompts)}, and, when depth $>1$ is allowed,
\code{rlm\_query(question, data)}, which starts a nested RLM on \code{data}. The
run ends when the model writes \code{FINAL(\ldots)} or \code{FINAL\_VAR(name)},
or sets \code{answer["ready"] = True}. A final that arrives in the same reply as
unexecuted code, or reads like a plan, is rejected once with a note. Thinking
(reasoning) text is logged but never fed back into the history. An optional
tool-calling protocol replaces fences and \code{FINAL} text with
\code{execute\_python} and \code{final\_answer} functions.

\subsection{Limits, compaction and forced finish}
Three classes of limit bound a run, and every one of them ends in an answer rather
than an exception.

\begin{itemize}[leftmargin=*]
\item \textbf{Sub-call caps} per code block and per run are enforced in code. A call
  past a cap raises an error inside the REPL, which the model sees.
\item \textbf{Context budget.} Every root request must fit under
  $\text{context\_tokens} - \text{max\_tokens}_{\text{root}} -
  0.1\,\text{context\_tokens}$ estimated tokens (never less than a quarter of the
  window), estimating one token per digit and 3.5 characters per token otherwise.
  Before each turn, REPL outputs older than the last four are stubbed and, if
  needed, summarised with one sub-call; if the request still does not fit, the
  remaining outputs are shrunk largest-first around a stub that points to the full
  text in a REPL \code{history} variable; if it still cannot fit, the request is not
  sent and the run takes the forced finish.
\item \textbf{Turns, time and errors.} When the turn cap, the wall-clock limit
  (\code{max\_timeout}) or the consecutive-error limit is reached, a single
  \emph{forced finish} call shows the model the current answer dictionary, the
  values of common answer variables and the end of the last REPL output, and asks
  for \code{FINAL} or \code{FINAL\_VAR}. Fallbacks, in order, are the answer
  dictionary, an answer variable, and the reply's prose with code and thinking
  removed; code is never returned as an answer. Every request is capped by the time
  left, so a run ends within \code{max\_timeout} plus 60 seconds plus about one
  second.
\end{itemize}

Several of these behaviours are fixes for defects found by live runs: a forced
finish that returned a stale REPL variable or a code block as the answer; an
error limit that discarded a run's work; a time limit that could be overrun by a
turn; and compaction that did not reserve room for the reply. Each is documented
in the repository's issue tracker \citep{reclamo}.

\subsection{Local serving: per-role endpoints and concurrency}
The root (planner) and sub (reader) roles may use different servers, keys and
concurrency limits; each distinct endpoint gets its own client and semaphore, so
sub-calls to a slow reader are not queued behind planner turns on another host.
\code{llm\_query\_batched} runs up to the reader's concurrency at once. On our main
server that is two parallel slots.

\subsection{Compatibility with the paper's RLM-Qwen3-8B}
The released RLM-Qwen3-8B model card states that the model ``was trained on
trajectories produced using a fixed system prompt'' and ``assumes the
environment/scaffold from our RLM repo'' \citep{rlmqwen8b}. The
\code{upstream-rlm-v0} planner style reproduces that scaffold: the paper's
Qwen3-8B/32K system prompt \citep[App.~C.1]{zhang2025rlm}, the context metadata
sent as an assistant message, the per-turn user prompt, \code{repl} fences, REPL
output formatting cut at 20,000 characters, and line-start \code{FINAL} parsing, as
in upstream \code{rlm} v1.0.0 \citep{rlmrepo}. Our sandbox, sub-call caps,
compaction and forced finish stay on. The exact training prompt is not public, so
parts of this reconstruction are approximate; the repository documents each
choice.

\subsection{Two prompt and limit configurations}
\label{sec:versions}
Table~\ref{tab:versions} lists the two configurations compared in this study.
Version 0.1 copied the paper's Qwen-specific warning against over-calling and
imposed conservative limits. Version 0.2 follows the paper's main (GPT-5) prompt
on delegation, adds the upstream ``orchestrator, not solver'' addendum, derives the
advertised sub-call size from the reader's window, and lifts limits we had imposed
ourselves. Both are selectable, so v0.1 results remain reproducible.

\begin{table}[t]
\centering\small
\caption{The two harness configurations. Sources are the RLM paper's appendices
\citep{zhang2025rlm} and the upstream code \citep{rlmrepo}. Sub-call size on
\code{pluto} assumes a 32,768-token reader window.}
\label{tab:versions}
\begin{tabularx}{\linewidth}{lXX}
\toprule
Setting & v0.1 & v0.2 \\
\midrule
Stance on sub-calls & ``sub-calls are expensive \ldots{} never 1000'' & ``strongly encouraged''; caps stated as facts \\
Planning guidance & none & orchestrator addendum (when to delegate and when not) \\
Advertised sub-call size & 12,000 characters & $(\text{window} - \text{max\_tokens}_{\text{sub}}) \times 0.85 \times 3$ characters \\
Root turns & 20 & 30 \\
REPL output shown & 2,000 characters & 20,000 characters \\
Root reply cap & 4,096 tokens & 8,192 tokens with thinking \\
Sub-call reply cap & 2,048 tokens & 4,096 tokens \\
Sub-calls per run / per block & 64 / 24 & 256 / 64 \\
Request timeout / run limit & 300 s / 1,800 s & 900 s / 3,600 s \\
\bottomrule
\end{tabularx}
\end{table}

\subsection{Security}
Generated code runs in a separate \code{python -I} process with a scrubbed
environment (no API keys), private protocol descriptors, and a per-execution
timeout; an optional Docker sandbox adds \code{--network none}, a read-only root
filesystem and a non-root user. The API key never reaches the REPL or the logs.

\subsection{Logging}
Every run writes a JSONL trajectory: metadata (never the key), each turn's content,
reasoning, code, outputs, notes and the serving endpoint and model, each sub-call's
sizes and latency, compaction and overflow events, and how the run ended. An
optional file records each root turn in the supervised-fine-tuning format the paper
used for RLM-Qwen3-8B.

% ===========================================================================
\section{An independently authored evaluation set}
\label{sec:tasks}

\paragraph{Why.} The harness, its prompts and its first tests were all written with
the help of Claude models (Section~\ref{sec:disclosure}). On those self-authored
tests (a million-line needle search, a synthetic ticket-classification task in the
style of OOLONG, and two-hop questions over a long digest) the harness looked
nearly perfect. Tests written by the harness's authors are likely to be shaped to
it, so we commissioned tests from models that had no part in building it.

\paragraph{How.} On 2026-10-07 a one-page brief went to eight non-Anthropic models
in the author's multi-model review harness \citep{flatline}, one lane at a time. Each
lane saw only the brief: no repository, no harness prompt and no other lane's
answer. The brief asked for one deterministic, standard-library Python generator
with \code{generate(seed, size)} and \code{score(answer, truth)}, producing
roughly 60,000, 300,000 and 1,200,000 characters of natural-looking text whose
answer requires breadth across the document and resists keyword search. A lane
running a Qwen model was excluded because the system under test is also a Qwen-branded model.
Table~\ref{tab:tasks} lists the eight tasks. The generators are kept byte-for-byte as
delivered; every repair lives in a separate file with its provenance, and every
version is pinned by SHA-256 in a manifest that CI checks.

\begin{table}[t]
\centering\small
\caption{The eight tasks. Each was written by a different non-Anthropic model,
named here by its lane in the review harness.}
\label{tab:tasks}
\begin{tabularx}{\linewidth}{llX}
\toprule
Lane & Author model & Question \\
\midrule
GLaDOS & xAI Grok 4.6 & Net authorized amount and certifying member for the legal successor of a renamed project, under bylaws and an SOP \\
SHODAN & OpenAI GPT-6 Astra & Earned freight credit per office after signed corrections, custody findings and agreement terms \\
TheDixieFlatline & Google Gemini 3.1 Pro & Final room of an asset through renames, transfers and corrections \\
Cerebex & Z.ai GLM-5.3 Flash & Travel total after amendments and reversals \\
Neuromancer & DeepSeek V4 Flash & March travel reimbursed after adjustments and denials \\
SELMA & NVIDIA Nemotron 3 Super & Final owner of a review after reassignments in an email thread \\
MasterControl & Mistral Medium 3.1 & The one employee flagged for two audit issues, and their total \\
Multivac & Nous Hermes 4 405B & Status and description of a requirement after merges and splits \\
\bottomrule
\end{tabularx}
\end{table}

\paragraph{Validation and repair.} Each generator was first checked in a
network-less container: its own self-test, the requested sizes, determinism within
and across processes, seed sensitivity, and that the true answer scores 1 and the
empty answer scores below 1. Five originals failed (a Python-version syntax
error, a degenerate small size, dependence on hash randomisation, a scorer that
split numbers, and contexts far below the requested size). Each failure went back
to its author model; in parallel, MiniMax-M3, another non-Anthropic model, repaired
all five. The author's fix was used when it passed, preserving independent
authorship. Later rounds found what these checks could not:

\begin{itemize}[leftmargin=*]
\item \textbf{Round 2} (found by the first evaluation run): one key described a
  different item from the one asked about, because question and answer drew
  separate random choices; one task never stated the starting state, so some
  answers were underivable; one task's ``unassigned'' event produced no text; and
  one scorer rejected correct answers written as sentences.
\item \textbf{Round 3} (an independent audit that solved every task from the text
  before opening any key, plus 1,701 scorer probes): one key filtered on attributes
  never stated in the text; one task's visible dates contradicted its event order;
  one correction email was never emitted; and in all eight tasks, correct answers in
  ordinary formats were marked down while hedged answers naming two candidates got
  full credit (251 flags).
\item \textbf{Round 3.1} (found by the pilot): one scorer gave 0.2 to a correct
  answer in an \code{ID: description --- Status: X} shape. The fix touched only
  \code{score()}; contexts, questions and keys stayed byte-identical. The scorer
  battery grew to 2,111 probes with zero flags and runs in CI.
\end{itemize}

Residual properties remain and are reported, not hidden: one task's key is
``Unassigned'' on about 62\% of seeds, and another's key is zero on about 26\%.

% ===========================================================================
\section{Evaluation design}
\label{sec:design}

\subsection{System under test}
The main model is distributed as Qwen3.8-Flash-Next, in an ``abliterated''
(refusal-removed) variant by Huihui.ai and a UD-Q4\_K\_XL quantization by Unsloth.
Despite the name, its GGUF metadata declares the architecture \code{qwen4exp}, ``A
Preview of the Qwen4 Architecture'': a mixture-of-experts model with 512 experts (10
active), 48 layers and a 262,144-token native context. It is therefore not a Qwen3
model, and any comparison with the RLM paper's Qwen3 results is cross-architecture.
It is served by the Strata inference server on \code{pluto},
a workstation with an NVIDIA V100 (32\,GB), a P100 (16\,GB) and an RTX 3060
(12\,GB, on an external PCIe dock). Since 2026-10-08 Strata runs with an int8 KV
cache, a 65,536-token resident cache and two parallel slots; it accepts prompts
beyond the resident cache at reduced prefill speed. The same model serves as
planner and reader. Root turns use thinking; sub-calls do not. Sampling follows the
Qwen3 model-card presets (root: temperature 0.6, top-$p$ 0.95, top-$k$ 20, min-$p$ 0;
sub-calls: 0.7, 0.8, 20, presence penalty 1.0), not the defaults embedded in the GGUF
(temperature 1.0, top-$p$ 0.95, top-$k$ 20). The full grid runs on the presets first.
The same grid is then repeated with the embedded defaults, and its results are kept
separate (H5).

\subsection{Arms}
All arms answer the same questions (Table~\ref{tab:arms}). ``32K'' is 32,768
tokens and ``128K'' is 131,072 tokens. The plain baseline is a single call with the
whole document in the prompt, thinking as the harness's root role sets it, and a
16,384-token output cap; a document that does not fit the usable window is recorded
as \emph{does not fit} and scores 0. There is no selection between attempts.

\begin{table}[t]
\centering\small
\caption{Arms. The plain arms differ only in window; the harness arms differ from
each other in configuration (Table~\ref{tab:versions}) and, for v0.2, window.}
\label{tab:arms}
\begin{tabularx}{\linewidth}{lX}
\toprule
Arm & Definition \\
\midrule
Plain 32K & One call, 32,768-token window (usable 24,576 after the reply reserve) \\
Plain 128K & One call, 131,072-token window, 1,800 s request timeout \\
Harness v0.1 & v0.1 prompt and limits, 32,768-token root window \\
Harness v0.2 32K & v0.2 prompt and limits, 32,768-token root window (pilot only) \\
Harness v0.2 128K & v0.2 prompt and limits, 131,072-token root window \\
\bottomrule
\end{tabularx}
\end{table}

\subsection{Questions and sample size}
Each cell is one task, size and seed. Small, medium and large contexts are about
60K, 300K and 1.2M characters; in the pilot, small documents were 16K--27K estimated tokens and medium ones
86K--148K; large documents run about 340K--610K (seed 0). Seeds 0 and
1 were used during development and are excluded. Seeds 2--3 (16 per cell) form the
pilot reported in Section~\ref{sec:lessons}.

\paragraph{Power.} McNemar's test uses only discordant pairs. In the pilot, Harness
v0.1 and Plain 128K disagreed on 9 of 30 attempted pairs, a discordance rate
$q \approx 0.30$. With $n \approx (z_{0.975} + z_{0.80})^2\, q / \delta^2$, detecting a
10-point difference $\delta$ at $\alpha = 0.05$ and 80\% power needs about 240 pairs
(about 320 with the interim looks below), and a 5-point difference about 940. A grid of
64 per cell detects only differences of roughly 25 points or more. The plan therefore
extends the grid in rolling batches toward 320 pairs (seeds 2--41) for the primary
comparison at small and medium; the first stage, seeds 2--9 (64 per cell, plus 64
large questions for the harness arms, which no plain arm can fit), is running.
5-point effects are out of reach on this hardware.

\paragraph{Pairing, clustering and replicates.} Every harness question is also answered
by plain; plain-only seeds may estimate baseline accuracy but never count toward paired
$n$. Variance is dominated by the task generator, so tasks are treated as a random effect:
we report per-task effects, a task-weighted average and a cluster bootstrap by task, and
note that generalising beyond these eight tasks needs more generators (the plan targets at
least 16). About 10\% of runs are same-question repeats, used to measure sampling noise and
never added to $n$. Runs of other serving stacks or model variants are separate
experiments, never pooled with \code{pluto}'s.

\subsection{Metrics and statistics}
Following the test plan in the repository (\code{evals/TEST-PLAN.md}), each row
records how the run ended (answered, forced finish after a limit, does not fit, or
not run because of a provider or transport failure, which is never scored). We
report completion separately from accuracy. Accuracy is the exact-match rate
(score 1.0 under the task's own scorer) with a Wilson 95\% interval, and the mean
score with a bootstrap 95\% interval. Paired comparisons use McNemar's exact test on
discordant cells and a paired bootstrap of the mean difference; a difference is
reported as real only if the test says so. Cost is tokens in and out including
sub-calls; latency is wall-clock time, reported as median and 95th percentile.
Every number names its code commit, task-set version, grader version, seeds and
profile. All tables are generated by \code{paper/analysis/make\_tables.py} from the
row files; none is typed by hand.

\subsection{Hypotheses}
\label{sec:hypotheses}
The test plan requires each phase's hypothesis, arms, sample size, metric, test and
stopping rule to be written down before the phase runs. The hypotheses below were
frozen in \code{evals/TEST-PLAN.md} (section 1a, commit \code{692a62a}) on 2026-10-09,
before any seed-4--9 score was examined. Those runs had started earlier that day, so the
freeze precedes looking at their results, though not their collection.

\begin{description}[style=unboxed,leftmargin=1.5em,font=\normalfont\bfseries]
\item[H1 (primary; harness vs plain, within the window).] At small and at medium
  size, exact-match accuracy of Harness v0.1 differs from Plain 128K on paired
  questions. Two-sided exact McNemar test with three looks at about 107, 213 and 320
  pairs, using O'Brien--Fleming boundaries ($|z| > 3.471, 2.454, 2.004$; corrected 2026-10-10 before any look, see \code{evals/TEST-PLAN.md} section 1b). At 64 pairs
  (seeds 2--9) no look is reached, so that grid reports estimates with intervals and no
  test of H1.
\item[H2 (secondary; v0.2 vs v0.1).] Harness v0.2 128K and Harness v0.1 differ in
  exact-match accuracy at small and at medium size (two-sided McNemar).
\item[H3 (secondary; delegation).] Harness v0.2 makes more sub-calls per question than v0.1
  (one-sided paired Wilcoxon signed-rank on sub-calls per question).
\item[H4 (beyond the window).] At large size, where no plain arm fits, both
  harness arms report completion and exact-match accuracy with intervals; no
  comparison with plain is possible. This is an estimate, not a test.
\item[H5 (secondary; sampling).] For each arm, accuracy with the model's embedded sampling
  defaults differs from accuracy with the Qwen3 presets on the same questions (two-sided
  exact McNemar per arm). H5 also asks whether the directions of H1 and H2 hold under
  both settings. The two grids are never pooled.
\end{description}

Ablations that attribute any H2 difference to the prompt, the turn cap or the REPL
output limit are a separate, later phase (Section~\ref{sec:future}).

% ===========================================================================
\section{Results}
\label{sec:results}

\PENDING{This section is filled by \code{make\_tables.py --stage main} once the
main grid reaches each pre-registered look: up to 320 paired questions for the primary
comparison at small and medium, starting with seeds 2--9 (64 per cell), and 64 large
questions per harness arm. No result below may be read as a finding until then.}

\IfFileExists{tables/main_accuracy.tex}{\input{tables/main_accuracy.tex}}{%
\pendingtable{Main grid: accuracy per arm and size, Wilson and bootstrap
intervals.}{main\_accuracy.tex, generated by \code{make\_tables.py --stage main}}}

\IfFileExists{tables/main_paired.tex}{\input{tables/main_paired.tex}}{%
\pendingtable{Main grid: paired McNemar tests for H1 and H2.}{main\_paired.tex}}

\IfFileExists{tables/main_cost.tex}{\input{tables/main_cost.tex}}{%
\pendingtable{Main grid: completion, delegation, tokens and latency.}{main\_cost.tex}}

\paragraph{H1.} \PENDING{outcome of H1 at small and at medium, with $p$-values.}

\paragraph{H2.} \PENDING{outcome of H2.}

\paragraph{H3.} \PENDING{outcome of H3.}

\paragraph{H4.} \PENDING{large-size completion and accuracy per harness arm.}

\paragraph{Breakdowns.} \PENDING{per-task and per-author-model accuracy, the
task-weighted average and the cluster bootstrap by task; replicate disagreement rate;
failure taxonomy (wrong evidence, misread, bad combination, early commit, turn cap,
format, crash).} \TODO{Add per-task and cluster-bootstrap tables to
\code{make\_tables.py}.}

% ===========================================================================
\section{Pilot and lessons learned (preliminary)}
\label{sec:lessons}

Everything in this section is preliminary. It comes from development runs and a
16-question-per-cell pilot, and it is reported because the lessons shaped the
design above, not because it supports any claim about the method.

\subsection{Lesson 1: the baseline was selected on the answer key}
Our first evaluations made two plain calls per question (thinking on and thinking
off), scored both against the key, and kept the better one, while the harness got
one attempt. That is an oracle. On the 16 small questions of a run on the round-2
task set, plain was exact on 12 as best-of-2, but on 9 with thinking on alone and on
8 with thinking off alone. The harness was exact on 6. The headline ``plain beats
the harness below the window'' from those runs was therefore partly an artefact of
selection. Since then plain is one call with no selection, and the two variants are
separate arms when needed.

\subsection{Lesson 2: graders moved results, and not evenly}
The first two evaluation runs used tasks that later audits found defective
(Section~\ref{sec:tasks}). Between the first run and a re-run after the round-2
fixes, harness exact answers on the four repaired tasks went from 6 of 20 to 12 of
20, and on the four unchanged tasks from 5 of 20 to 6 of 20: the apparent
improvement came from repairing tasks, not from changing the harness. In the pilot
below, the round-3.1 scorer fix changed four Plain 128K cells from 0.2 or 0.3 to 1.0,
one cell in each v0.2 harness arm and in Plain 32K, and none in Harness v0.1. A grader
defect thus cost the baseline more than the harness. Four Plain 128K answers could
not be re-scored because the row files stored only their first 2,000 characters;
they keep their original scores.

\subsection{Lesson 3: a prompt fix for over-delegation suppressed delegation}
The paper's Qwen-specific prompts add a warning against sub-calls because
Qwen3-Coder over-delegated \citep[App.~C]{zhang2025rlm}. We copied that warning
into v0.1. Our model has the opposite tendency: on the round-2 re-run only 7 of 40
harness rows made any sub-call, and 15 of 40 used all 20 turns, mostly reading
slices of the document itself and writing ad hoc parsers. Version 0.2 removes the
warning and lifts our limits. In the pilot it delegated far more (sub-calls on
14 of 16 small questions at 128K, against 3 of 16 for v0.1; Table~\ref{tab:pilot-cost})
without a detectable change in accuracy (Table~\ref{tab:pilot-paired}).

\subsection{The pilot grid}
The pilot ran the five arms of Table~\ref{tab:arms} on seeds 2--3 of the round-3 task
set, graded with the round-3.1 scorers (Tables~\ref{tab:pilot-accuracy},
\ref{tab:pilot-cost} and \ref{tab:pilot-paired}). Harness v0.1 was exact on
\pVoneS{} small and \pVoneM{} medium questions; Plain 128K on \pPlainLongS{} and
\pPlainLongM{}. Pooled over both sizes the discordant cells split
\pBCVoneVsPlainLong{} (McNemar $p = \pPVoneVsPlainLong$): \emph{no detectable
difference at this sample size}. Harness v0.2 128K against Plain 128K split
\pBCVtwoVsPlainLong{} ($p = \pPVtwoVsPlainLong$), and against Harness v0.1
\pBCVoneVsVtwoLong{} ($p = \pPVoneVsVtwoLong$). The only significant pilot
contrasts are those where the plain prompt does not fit at all (Plain 32K at medium
size), which are true by construction. A few bootstrap intervals on paired mean
differences exclude zero (Table~\ref{tab:pilot-paired}); at $n = 16$, uncorrected
for the number of comparisons, we treat none of them as evidence. Sixteen questions per cell cannot
distinguish differences of a few questions; the pilot's discordance rate is what sized
the main study (Section~\ref{sec:design}).

\begin{table}[t]
\centering\small
\caption{Accuracy per arm and size. \textbf{Preliminary pilot}: seeds 2, 3, $n{=}16$ cells per arm and size, round-3.1 grader. Not a test of any hypothesis. Exact = score 1.0 under the task's own \texttt{score()}; a plain prompt that does not fit the window (DNF) scores 0. Intervals: Wilson 95\% for exact, bootstrap 95\% for the mean.}
\label{tab:pilot-accuracy}
\begin{tabular}{llrlrlr}
\toprule
Size & Arm & Exact & 95\% CI & Mean & 95\% CI & DNF \\
\midrule
Small & Plain 32K & 9/16 & [0.33, 0.77] & 0.60 & [0.37, 0.83] & 2 \\
 & Plain 128K & 8/16 & [0.28, 0.72] & 0.56 & [0.34, 0.78] & 0 \\
 & Harness v0.1 & 11/16 & [0.44, 0.86] & 0.71 & [0.49, 0.90] & 0 \\
 & Harness v0.2 32K & 10/16 & [0.39, 0.82] & 0.68 & [0.45, 0.88] & 0 \\
 & Harness v0.2 128K & 10/16 & [0.39, 0.82] & 0.69 & [0.48, 0.89] & 0 \\
\midrule
Medium & Plain 32K & 0/16 & [0.00, 0.19] & 0.00 & [0.00, 0.00] & 16 \\
 & Plain 128K & 7/16 & [0.23, 0.67] & 0.47 & [0.24, 0.71] & 2 \\
 & Harness v0.1 & 9/16 & [0.33, 0.77] & 0.62 & [0.39, 0.83] & 0 \\
 & Harness v0.2 32K & 7/16 & [0.23, 0.67] & 0.49 & [0.27, 0.72] & 0 \\
 & Harness v0.2 128K & 8/16 & [0.28, 0.72] & 0.57 & [0.34, 0.78] & 0 \\
\bottomrule
\end{tabular}
\end{table}

\begin{table}[t]
\centering\small
\setlength{\tabcolsep}{4pt}
\caption{Paired comparisons on identical cells. \textbf{Preliminary pilot}: seeds 2, 3, $n{=}16$ cells per arm and size, round-3.1 grader. Not a test of any hypothesis. $b$ = cells only A got exact, $c$ = cells only B got exact; $p$ = McNemar's exact two-sided test on $(b, c)$. $\Delta$ = paired mean-score difference A$-$B with a bootstrap 95\% interval. Uncorrected for multiple comparisons.}
\label{tab:pilot-paired}
\begin{tabular}{lllrrrrl}
\toprule
Size & A & B & A ex. & B ex. & $b$/$c$ & $p$ & $\Delta$ mean [95\% CI] \\
\midrule
Small & Harness v0.1 & Plain 128K & 11/16 & 8/16 & 5/2 & 0.45 & $+0.15$ [-0.14, 0.42] \\
 & Harness v0.1 & Plain 32K & 11/16 & 9/16 & 5/3 & 0.73 & $+0.11$ [-0.23, 0.42] \\
 & Harness v0.2 128K & Plain 128K & 10/16 & 8/16 & 2/0 & 0.50 & $+0.13$ [0.01, 0.29] \\
 & Harness v0.1 & Harness v0.2 128K & 11/16 & 10/16 & 3/2 & 1.00 & $+0.02$ [-0.23, 0.25] \\
 & Harness v0.1 & Harness v0.2 32K & 11/16 & 10/16 & 1/0 & 1.00 & $+0.04$ [0.00, 0.11] \\
\midrule
Medium & Harness v0.1 & Plain 128K & 9/16 & 7/16 & 2/0 & 0.50 & $+0.14$ [0.03, 0.30] \\
 & Harness v0.1 & Plain 32K & 9/16 & 0/16 & 9/0 & 0.004 & $+0.62$ [0.40, 0.83] \\
 & Harness v0.2 128K & Plain 128K & 8/16 & 7/16 & 3/2 & 1.00 & $+0.09$ [-0.14, 0.33] \\
 & Harness v0.1 & Harness v0.2 128K & 9/16 & 8/16 & 2/1 & 1.00 & $+0.05$ [-0.11, 0.24] \\
 & Harness v0.1 & Harness v0.2 32K & 9/16 & 7/16 & 3/1 & 0.62 & $+0.12$ [-0.07, 0.34] \\
\midrule
Pooled & Harness v0.1 & Plain 128K & 20/32 & 15/32 & 7/2 & 0.18 & $+0.15$ [-0.01, 0.31] \\
 & Harness v0.1 & Plain 32K & 20/32 & 9/32 & 14/3 & 0.01 & $+0.37$ [0.14, 0.57] \\
 & Harness v0.2 128K & Plain 128K & 18/32 & 15/32 & 5/2 & 0.45 & $+0.11$ [-0.02, 0.25] \\
 & Harness v0.1 & Harness v0.2 128K & 20/32 & 18/32 & 5/3 & 0.73 & $+0.03$ [-0.11, 0.18] \\
 & Harness v0.1 & Harness v0.2 32K & 20/32 & 17/32 & 4/1 & 0.38 & $+0.08$ [-0.02, 0.20] \\
\bottomrule
\end{tabular}
\end{table}

\begin{table}[t]
\centering\small
\setlength{\tabcolsep}{4pt}
\caption{Completion, delegation, cost and latency over attempted rows. \textbf{Preliminary pilot}: seeds 2, 3, $n{=}16$ cells per arm and size, round-3.1 grader. Not a test of any hypothesis. Ans.\ = the model finished on its own (\texttt{FINAL}, \texttt{FINAL\_VAR}, the answer dict, or a plain reply that stopped normally); Lim.\ = turn cap, time limit, error limit, context overflow or output-length cap (the harness then takes a forced finish). Deleg.\ = rows with at least one sub-call; Sub-calls and Turns are means per attempted row. Tokens include sub-calls. Seconds are wall clock on shared, serial hardware.}
\label{tab:pilot-cost}
\begin{tabular}{llrrrrrrr}
\toprule
Size & Arm & Ans. & Lim. & Deleg. & Sub-calls & Turns & Med.\ tokens & Med./p95 s \\
\midrule
Small & Plain 32K & 14/14 & 0/14 & -- & -- & -- & 17{,}870 & 113/226 \\
 & Plain 128K & 16/16 & 0/16 & -- & -- & -- & 18{,}873 & 134/263 \\
 & Harness v0.1 & 12/16 & 4/16 & 3/16 & 0.75 & 13.50 & 55{,}137 & 128/515 \\
 & Harness v0.2 32K & 14/16 & 2/16 & 10/16 & 1.62 & 12.56 & 43{,}277 & 206/1{,}201 \\
 & Harness v0.2 128K & 16/16 & 0/16 & 14/16 & 2.44 & 12.25 & 67{,}688 & 252/1{,}856 \\
\midrule
Medium & Plain 32K & 0/0 & 0/0 & -- & -- & -- & -- & --/-- \\
 & Plain 128K & 12/14 & 2/14 & -- & -- & -- & 88{,}699 & 449/607 \\
 & Harness v0.1 & 10/16 & 6/16 & 4/16 & 0.38 & 14.75 & 103{,}628 & 236/459 \\
 & Harness v0.2 32K & 14/16 & 2/16 & 7/16 & 1.44 & 16.50 & 100{,}321 & 371/2{,}739 \\
 & Harness v0.2 128K & 15/16 & 1/16 & 10/16 & 2.06 & 16.75 & 132{,}842 & 374/3{,}245 \\
\bottomrule
\end{tabular}
\end{table}


\subsection{Side runs with other planners}
Three side runs, on the round-3 task set with v0.1 settings, are summarised in
Table~\ref{tab:pilot-side}. Their plain rows predate Lesson 1 and are best-of-2.

\paragraph{The paper's RLM-Qwen3-8B as planner.} With the compatibility mode, the
released checkpoint (a community GGUF conversion on a local llama.cpp server) planned
and Flash-Next on \code{pluto} read. It was exact on 1 of 16 small, 4 of 16
medium and 0 of 8 large questions. It committed early (median 2 turns), and 14 of
40 rows ended at the time limit.
\emph{These conditions differ materially from the paper's}: the planner window was
16,384 tokens rather than 32K; the reader was a different model on slower, shared
hardware rather than Qwen3-8B; our sub-call caps and error messages are ones the
checkpoint never saw in training; the small rows and the medium and large rows ran at
different harness commits; and the reader's server configuration changed during the
large rows. Nor does it test the paper's actual 8B claim, that fine-tuning helps:
there was no base Qwen3-8B arm. The result says nothing about RLM-Qwen3-8B under its
native conditions, which remain to be tested (Section~\ref{sec:future}).

\paragraph{Poolside Laguna S 2.1.} With this hosted model the harness was exact on 5
of 16 small, 2 of 16 medium and 3 of 8 large questions, and 28 of 40 runs hit the
20-turn cap. Laguna writes its native tool-call syntax
(\code{<tool\_call>repl \ldots</tool\_call>}) as message text, which the fence
parser originally ignored, turning those turns into no-ops; it also sometimes
followed a call with REPL output it had invented. The parser now accepts the form
and discards the rest of a reply after the first such call. Re-running 24
cells after the fix moved exact answers from 5 to 6: the fix restored the protocol but did
not detectably change accuracy. A provider quota exhaustion during this run produced
rows recorded as scores of zero; they were quarantined, and the runner now must
record provider failures as not run.

\paragraph{Llama 3.3 70B, and Python versus Rust.} A separately developed Rust
port of the harness scored poorly with Llama 3.3 70B. To separate the language of the
harness from the model, we ran the Python harness with the same model, provider and
sampling on the same GLaDOS cells: neither the harness nor plain was exact on any of
the four cells (mean scores 0.10 for both modes). \TODO{Insert the Rust port's
per-cell scores from its repository and confirm they match; the claim of equal
scores comes from the session log and is not yet verified against the Rust run
files.} The working conclusion is that the model, not the harness's language, was
the bottleneck.

\IfFileExists{tables/pilot_side.tex}{\begin{table}[t]
\centering\small
\caption{\textbf{Preliminary} pilot side runs on the round-3 task set (seeds 0--1 at small and medium, seed 0 at large; the Llama run is GLaDOS only), with harness v0.1 settings. Their plain rows still used the best-of-2 rule of Section~\ref{sec:lessons}, so plain is an upper bound. Exact / cells. Observations, not tests.}
\label{tab:pilot-side}
\begin{tabular}{llrr}
\toprule
Setup & Size & Harness exact & Plain exact (best-of-2) \\
\midrule
RLM-Qwen3-8B planner & small & 1/16 & -- \\
 & medium & 4/16 & -- \\
 & large & 0/8 & -- \\
\midrule
Laguna S 2.1 & small & 5/16 & 4/16 \\
 & medium & 2/16 & 2/16 (2 DNF) \\
 & large & 3/8 & -- \\
\midrule
Llama 3.3 70B (GLaDOS only) & small & 0/2 & 0/2 \\
 & medium & 0/2 & 0/2 (2 DNF) \\
\midrule
Laguna, parser-fix reruns & 24 cells & 5/24 before, 6/24 after & -- \\
\bottomrule
\end{tabular}
\end{table}
}{%
\pendingtable{Pilot side runs.}{pilot\_side.tex, generated with \code{--side}}}

\subsection{Practicality}
Wall-clock time, not accuracy, was the binding constraint throughout. Sub-calls run
one or two at a time on a serial server; a plain 128K call on a medium document spends
several minutes reading its prompt (median \PENDING{main-grid median} s); and the
95th-percentile harness run under v0.2 took close to an hour in the pilot
(Table~\ref{tab:pilot-cost}). Hardware failures cost more time than any single
experiment: on 2026-10-09 the RTX 3060 fell off the PCIe bus (NVIDIA Xid 79) for
about 2.5 hours, turning 14 harness rows into errors that the runner then recorded
as zero scores; a network interface failed later the same morning. All affected
rows were removed (the originals are kept) and re-run. This matches the paper's own
account. Its limitations section says that ``future strategies involving
asynchronous sub-calls and sandboxed REPLs can potentially significantly reduce the
runtime and inference cost of RLMs'' \citep[\S7]{zhang2025rlm}, and its negative
results say: ``RLMs without asynchronous LM calls are slow. We implemented all
sub-LM queries naively as blocking / sequential calls, which caused our RLM
experiments to be slow, especially compared to just the base model''
\citep[App.~B]{zhang2025rlm}.

\subsection{Process errors}
For completeness, errors in how the study was run, each fixed and recorded in the
repository: a pull request merged without waiting for CI broke the main branch on
one Python version; the first validator checked only that a key scores itself, and
so missed the answerability defects; merging a pull request moved the experiment's
checkout mid-run (a scorer-only change), crashing one arm, which was resumed;
output tokens were reserved twice in the 8B runs, producing context-overflow rows that
were discarded and re-run; and turn-by-turn trajectories of three early development
runs were lost when temporary worktrees were cleaned (their summary files survive).

% ===========================================================================
\section{Discussion}
\label{sec:discussion}

\PENDING{Interpretation of H1--H4. To be written after the main grid. It should
address: whether any harness advantage within the window survives fair measurement
at the pre-registered sample size; whether more delegation buys accuracy on these tasks; and the
cost--accuracy trade-off on serial hardware.}

Two points can be made now, because they concern design rather than results. First,
our tasks are mostly multi-hop chains whose deciding evidence is a handful of
records, which in the taxonomy of \citet{goldman2025retrieval} is low-dispersion
work. The RLM paper's largest gains are on information-dense aggregation such as
OOLONG and OOLONG-Pairs, and its earlier versions observe that ``the base LM
outperforms RLM in the small input context regime''
\citep[v2, \S4]{zhang2025rlm}. A null result on our tasks would therefore not contradict the paper.
Second, the pilot shows that measurement choices (baseline selection, grader
defects, a copied prompt patch) can move a comparison of this kind by several
questions in 16, which is the size of the effects under study.

% ===========================================================================
\section{Threats to validity}
\label{sec:threats}

\begin{itemize}[leftmargin=*]
\item \textbf{Synthetic data.} All tasks are generated corporate records written by
  other models. They may share conventions that favour or penalise code-based
  reading, and they are not real documents.
\item \textbf{One main model, of a different architecture from the paper's.} Results
  are for one quantized, abliterated open-weight model on one server. Its architecture
  (\code{qwen4exp}) differs from every Qwen model in the RLM paper, so within-model
  comparisons (harness vs plain, v0.1 vs v0.2) are valid but comparisons with the paper
  are cross-architecture. Abliteration and 4-bit quantization may change behaviour
  beyond refusals.
\item \textbf{Sampling and templates.} All runs used Qwen3 model-card sampling presets
  rather than the model's embedded defaults, and the thinking switch was assumed to
  behave on this architecture's chat template as on Qwen3's. Neither is verified yet.
\item \textbf{Small samples, few tasks.} The pilot has 16 questions per cell; the main
  study targets 320 pairs for the primary comparison but starts at 64. With eight task
  generators, effective sample size is limited by between-task variance. Run-to-run noise
  is not yet measured.
\item \textbf{Shared hardware.} \code{pluto} served other work during the study;
  latency figures reflect contention, two parallel slots, and a server configuration
  that changed on 2026-10-08.
\item \textbf{Code changes during runs.} Two arms crossed a scorer-only commit;
  side runs crossed harness commits. Every row records its commit, and pilot rows were
  re-graded with a single grader version.
\item \textbf{Outages and provider failures.} Rows lost to hardware or provider
  failures were removed and re-run, with originals kept; before that rule existed,
  such rows were scored as zero.
\item \textbf{Grader residuals.} The scorers pass a 2,111-probe battery but are
  still programmatic. Human spot-checks of 10\% of graded answers are planned, not
  done. Four plain answers in the pilot could not be re-graded.
\item \textbf{Author and assistant overlap.} The harness and its prompts were
  developed with Claude models, and Claude also orchestrated the evaluation. The
  task generators were deliberately sourced elsewhere, but the analysis pipeline was
  not independently reimplemented.
\end{itemize}

% ===========================================================================
\section{Future work}
\label{sec:future}

The test plan orders the remaining work as follows. (1) Finish the first stage of the main grid (seeds 4--9, then 64 large questions) and extend the primary comparison toward 320 pairs. (2) Validity checks: question-only and
evidence-deleted controls, counterfactual twins, independent hand reconstruction of
keys, a 10\% human grading audit, and recovery of answers reached but not submitted.
(3) Instrumentation of completion and of what each arm actually received. (4) A
one-factor-at-a-time decomposition of v0.1 versus v0.2 (prompt, turn cap, REPL
output limit), plus a check at matched total compute. (5) Oracle interventions that
give readers the correct evidence or replace their outputs, and a BM25 retrieval
baseline. (6) Coverage: unanswerable questions mixed 50/50, a size sweep from 16K to
1M tokens to locate any crossover, paraphrased questions, evidence position and
chunk-boundary placement, generator-convention sensitivity, aggregation and
long-output tasks, contradictions, real documents, and a reproduction of a paper
benchmark (OOLONG or BrowseComp-Plus). (7) A replication of the paper's fine-tuning claim: base Qwen3-8B against
RLM-Qwen3-8B as planner in the same harness and on the same questions, under native
conditions (32K window, 20,000-character REPL outputs, Qwen3-8B reader, with a
Flash-Next-reader variant as a separate arm), checked first on the paper's own tasks
against its published numbers. This is also the same-architecture anchor that the main
model cannot provide. (8) On rented GPUs: planner/reader pairings of different sizes.
Also planned: the repeat of the main grid with the model's embedded sampling defaults
(H5), at least eight
more task generators from outside models, and a frontier model as a ceiling. 

% ===========================================================================
\section*{Code and data availability}
The harness, the task generators with their full provenance (brief, verbatim
answers, every repair round and the audit), the runner and this paper's analysis
script are at \url{https://github.com/CryptoJones/ReCLamO-Harness}. Run row files and
trajectories are archived off-machine and are not yet in the repository.
\TODO{Publish the exp-v02 row files and trajectories (scrubbed of host addresses)
with the final version.}

\section*{AI assistance disclosure}
\label{sec:disclosure}
AI systems are not authors of this paper. Claude (Anthropic; Opus models)
orchestrated the project, analysed results, wrote code, and drafted text including
this manuscript and its analysis script. Claude Fable (Anthropic) wrote the v0.1
implementation of the harness. MiniMax-M3 repaired task generators and, together
with a roundtable of non-Anthropic models (GLM, Mistral, DeepSeek, Nemotron, GPT,
Gemini, Grok, Hermes and Qwen families), authored the independent test generators and
reviewed the evaluation design (Grok's lane failed to authenticate for the design
review), through the author's multi-model review harness
\citep{flatline}, also used in \citet{clark2026veto}. A Qwen-family lane reviewed the
design but did not author a test generator, because the system under test is also a Qwen-branded model. A separately developed Rust port of the harness, used only for the
cross-check in Section~\ref{sec:lessons}, was written by MiniMax-M3. Model outputs were checked against the repository, the run files and the cited
sources before use. The author directed the work, made the decisions recorded in the
repository, and takes responsibility for the content, including any errors.

\section*{Acknowledgments}
The author thanks Alex L.\ Zhang for encouragement: on seeing the harness he replied,
``This is awesome! You should do a write up of your findings.'' This draft is that
write-up, in progress. Any errors in reproducing the RLM method are the author's,
not the original authors'.

\section*{Declarations}
\noindent\textbf{Funding:} none. \quad \textbf{Competing interests:} none.

\bibliographystyle{plainnat}
\bibliography{refs}

\end{document}
